\section{Extension O: covariance infrastructure (Stage 8B.1, partial)}

This section documents only what is implemented in \texttt{src/fit\_covariance.py}. It is the first part of a larger plan; it does not complete the fit-covariance work.

\paragraph{Implemented.} (i) \texttt{MeasurementSet}: named values with units, a validated covariance and a provenance (\emph{measured}, \emph{theory}, \emph{recommended}, \emph{scenario}); the arrays are read-only. (ii) \texttt{validate\_covariance}: finite values, symmetry and positive-semidefiniteness with tolerances relative to the largest entry and eigenvalue. (iii) \texttt{numerical\_jacobian}: central differences at three step sizes with a stability check. (iv) \texttt{inverse\_map\_covariance}: propagation through an explicit map $\mathbf p=\mathbf p(\mathbf y)$. (v) \texttt{identifiability}: singular-value diagnostics of a scaled Jacobian.

\paragraph{Not implemented.} Weighted-fit covariance $(J^{T}C_y^{-1}J)^{-1}$, residual-variance covariance $s^2(J^{T}J)^{-1}$, a general Monte Carlo engine, helium numerical covariance, and any change to \texttt{fit\_parameters} or \texttt{fit\_to\_levels}. In particular, the unweighted \texttt{fit\_parameters} is not given a weighted covariance, and the constrained calcium model \texttt{fit\_to\_levels} receives no covariance at all: its residuals are model discrepancy, and noise is not estimated from them.

\paragraph{Propagation.} For $C_y$ the covariance of the inputs and $A=\partial\mathbf p/\partial\mathbf y$,
\begin{equation}
C_p = A\,C_y\,A^{T}.
\end{equation}
This is a first-order, local statement. It is valid for any positive-semidefinite $C_y$, including a singular one (for example a common shift of all inputs), so zero eigenvalues are not treated as invalid. The Jacobian is analytical when supplied, otherwise numerical; a numerical Jacobian whose three step sizes ($h/2$, $h$, $2h$) disagree by more than $10^{-5}$ of the column scale is refused and no covariance is returned. If the map raises an error at a perturbed point (outside its domain) a Jacobian error is raised instead of a number.

\paragraph{Validation behaviour.} A covariance must be finite, square, symmetric ($\max|C-C^{T}|\le 10^{-10}\max|C|$) and positive semidefinite ($\lambda_{\min}\ge-10^{-10}\lambda_{\max}$). Non-symmetric and negative-eigenvalue matrices are rejected. Correlations are never inferred, for instance from two numbers appearing in one publication: they are supplied explicitly or the covariance is diagonal. A later weighted fit needs $C_y^{-1}$; \texttt{require\_positive\_definite} and \texttt{whitening\_factor} reject a singular matrix with an explanatory error and apply neither regularisation nor a pseudoinverse.

\paragraph{Identifiability and scaling.} The diagnostic analyses $J_{\mathrm{scaled}}=L^{-1}J S_p$ for a forward Jacobian $J=\partial f/\partial p$ ($C_y=LL^{T}$ positive definite, $S_p$ the diagonal matrix of characteristic parameter scales), and $S_p^{-1}A\,\mathrm{diag}(\sigma_y)$ for the inverse-map Jacobian. It reports the singular values, the numerical rank and tolerance (default $\max(\text{shape})\,\epsilon\,s_{\max}$), the condition number $s_{\max}/s_{\min}$, the weak parameter combinations (singular vectors) and the largest parameter correlation, together with a status \emph{ok}, \emph{flagged} or \emph{refused}. The scaling convention is stored with every result, because a condition number is meaningless without it; the default scale is $|\hat p|$. The thresholds $10^{6}$ (flag) and $10^{8}$ (refuse) are configurable diagnostic conventions, not physical limits, and a large condition number does not by itself prove physical non-identifiability. For outputs that are derived from other outputs (calcium $\varepsilon^2$ next to its four parameters) a rank deficiency is flagged but the singular output covariance is returned.

\paragraph{Calcium regression.} The four-parameter calcium map, with an assumed independent scenario of 0.001~cm$^{-1}$ per level, reproduces \texttt{calcium\_fit\_sensitivity} for all five outputs: standard deviations agree to $2\times10^{-8}$ relative, correlation matrices to $1\times10^{-8}$, and the analytic derivatives of $\varepsilon^{2}$ are matched. The result is $\sigma(\varepsilon^2)=1.477\times10^{-7}$ and is labelled a sensitivity scenario. For the four parameters alone, the condition number is $3.7\times10^{4}$ (scaling $S_p=|\hat p|$ and inputs scaled by their $\sigma$) and the largest parameter correlation is 0.80. An explicitly correlated scenario (all pairwise level correlations 0.5) is checked against analytic expressions, and a common shift of all four levels (a rank-one covariance) leaves $\zeta_S$, $\zeta_A$, $\Delta$ and $\varepsilon^2$ unchanged and moves only $E_t$, as it must.

\paragraph{Scenario versus measured.} A \emph{scenario} covariance is an assumed one, for example that every calcium level is uncertain by 0.001~cm$^{-1}$ independently, chosen at about the scale of the tabulated precision; NIST publishes no uncertainties for these levels. Its propagated result shows how sensitive the output is to errors of that size, and is not an experimental uncertainty. A \emph{measured} covariance comes from a source that states the uncertainties and their correlations.

\paragraph{What the covariance does not establish.} It is conditional on the map and contains no model error (missing perturbers, omitted interactions) and no numerical error of the map itself. It is not a goodness-of-fit: an exact calibration with zero degrees of freedom has no residual to test, and agreement between a model and the levels used to calibrate it is not validation. Local linear propagation fails near the boundary of the map's domain or where the map curves on the scale of the input uncertainty. For calcium, $V^2=(a-E_{3P_1})(E_{1P_1}-a)\ge0$ is such a boundary: a displacement $a-E_{3P_1}=0.518$~cm$^{-1}$ is only a few standard deviations away once the assumed per-level uncertainty approaches $0.1$~cm$^{-1}$, and sampling then meets inputs for which no real $\zeta_A$ exists.

\paragraph{Deferred requirement for future Monte Carlo code.} A sampler must separate expected domain rejection (such as $V^2<0$) from numerical failure and from unexpected programming exceptions; it must report the number of rejected samples and the acceptance fraction; and statistics of the accepted samples are conditional on acceptance and must be labelled so. No such engine exists yet.
